
\chapter{Системийн шинжилгээ}
\label{Chapter2}

\section{Систем боловсруулах хэрэгцээ}
Онолын судалгаанаас харахад програмчлалын сургалтын цахим орчинд дараах үндсэн асуудлууд нийтлэг байна: хэрэглэгчийн мэдлэгийн ялгааг тооцохгүй ижил урсгал, ахицын өгөгдлийг бүрэн ашигладаггүй сул аналитик, багшгүй үед тайлбарын хоцрогдол, суралцах идэвхийг тогтвортой хадгалах хүндрэл. Иймээс шинэ систем нь зөвхөн хичээлийн агуулга байршуулдаг платформ бус, харин хэрэглэгчийг эхлэлээс нь удирдан дагуулдаг, агуулгыг шаталсан байдлаар санал болгодог, Хиймэл оюун-аар тайлбар дэмждэг ухаалаг веб систем байх шаардлагатай.

Энэхүү дипломын ажлын систем нь хэрэглэгчийн анхны түвшин тогтоохоос эхлээд дараагийн хичээлийг санал болгох, хичээл тутмын ойлголтыг Хиймэл оюун-аар бататгах, суралцах ахицыг тоон мэдээллээр харах, community хэлбэрээр харилцах, тоглоомчилсон контент ашиглах боломжийг нэг дор төвлөрүүлснээр дээрх асуудлыг шийдвэрлэхийг зорьсон.

\section{Системийн зорилго ба хэрэглээний орчин}
Системийн ерөнхий зорилго нь програмчлалын агуулгад чиглэсэн, хэрэглэгчийн түвшин болон ахицад нийцсэн шаталсан сургалтын платформ хөгжүүлэх юм. Хэрэглэгч системд бүртгүүлж, Түвшин тогтоох шалгалт өгсний дараа өөрт тохирсон түвшинд хамаарах content-ийг ашиглаж эхэлнэ. Хичээл бүр дээр тайлбар, хураангуй, асуулт, дасгал ажил, хэлэлцүүлэг зэрэг нэмэлт хэрэгсэл байгаа нь суралцах үр дүнг сайжруулахад чиглэнэ.

Энэ системийг дараах орчинд ашиглах боломжтой.
\begin{itemize}
\item Их, дээд сургуулийн хичээлийн нэмэлт дэмжих платформ
\item Програмчлал анхан шатнаас бие даан сурч буй хэрэглэгчдийн онлайн орчин
\item Сургалтын төвийн түвшин ялгасан холмиг сургалт платформ
\item Туршилтын хэлбэрийн Хиймэл оюуны тусламжит e-learning prototype
\end{itemize}

\section{Хэрэглэгчийн бүлэг ба тэдгээрийн хэрэгцээ}
Системийн үндсэн хэрэглэгч нь суралцагч боловч дүрийн түвшинд student, mentor, admin гэсэн ялгаа бий. Кодын түвшинд хэрэглэгчийн \texttt{role} болон \texttt{skill\_level} талбарууд хадгалагддаг. Үүний дагуу шаардлагыг дараах байдлаар задлан авч үзэв.

\subsection{Суралцагч хэрэглэгч}
Суралцагч нь системийн гол урсгалыг ашиглана. Түүнд дараах хэрэгцээ бий.
\begin{itemize}
\item Бүртгүүлэх, нэвтрэх, профайлаа удирдах
\item Анхны түвшнээ автоматаар тодорхойлуулах
\item Өөрийн түвшинд тохирсон course, хичээл харах
\item Хичээлийн урт агуулгыг товч хураангуй хэлбэрээр унших
\item Ойлголтоо асуулт болон дасгал ажил-аар шалгах
\item Долоо хоногийн идэвх, нийт оноо, курсын ахицаа харах
\item Түвшин ахих шалгалт-аар дараагийн түвшин рүү ахих
\item Бусад хэрэглэгчтэй community хэсэгт харилцах
\end{itemize}

\subsection{Mentor хэрэглэгч}
Mentor түвшин нь ахисан хэрэглэгч эсвэл тодорхой шалгуур хангасан суралцагчаас үүсч болно. Одоогийн кодод advanced түвшинд хүрч, тодорхой тооны хичээл дуусгасан тохиолдолд role автоматаар mentor руу шилжих логик байгаа. Mentor хэрэглэгч контент үүсгэх эсвэл удирдах түвшний эрх авах боломжтой байдлаар систем өргөтгөхөд бэлэн.

\subsection{Admin хэрэглэгч}
Admin нь системийн бүрэн эрхтэй хэрэглэгч бөгөөд backend admin хэсэг, өгөгдөл удирдлага, контентын засвар, орчны тохиргоонд оролцох боломжтой. Төслийн өнөөгийн хэрэгжилтэд контент бүтэц, хэрэглэгчийн мэдээлэл, landing page өгөгдөл, media файлуудыг удирдах суурь бүрдсэн.

\section{Бизнес процессын шинжилгээ}
Системийн үндсэн урсгал нь дараах дараалалтай.
\insertimagecustom{zurag-diagram/flowchart.png}{Системийн үндсэн бизнес процессын үйл ажиллагааны диаграм}{height=0.78\textheight,keepaspectratio}
\begin{enumerate}
\item Хэрэглэгч бүртгүүлж, нэвтэрнэ.
\item Хэрэв Түвшин тогтоох шалгалт өгөөгүй бол эхний түвшинг тодорхойлох шалгалт руу орно.
\item Түвшин тогтоох шалгалт-ийн үр дүнд үндэслэн skill level хадгалагдана.
\item Систем хэрэглэгчид тохирсон course, track, хичээл харуулна.
\item Хэрэглэгч хичээл үзэж, summary, асуулт, дадлага ажил, discussion ашиглана.
\item хичээл completion болон score нь progress мэдээлэлд бүртгэгдэнэ.
\item Нийт оноо, completed хичээл-ийн тоо өсөхийн хэрээр skill level болон role шинэчлэгдэнэ.
\item Хэрэглэгч Түвшин ахих шалгалт өгч дараагийн шатанд ахих боломжтой.
\item Progress summary хэсэгт долоо хоногийн идэвх, course бүрийн гүйцэтгэлээ хянаж, цааш сурах чиглэлээ тодорхойлно.
\end{enumerate}

Энэ урсгал нь статик сургалтын системээс ялгаатай нь хэрэглэгчийн үйлдэл бүрийг дараагийн алхамтай нь холбоно. Жишээлбэл хичээл дуусгалт нь зөвхөн ``дууссан'' төлөв хадгалахаас гадна нийт оноо, ур чадварын түвшин, санал болголт, курсын ахицад нөлөөлнө.

\section{Функциональ шаардлага}
Системийн функциональ шаардлагыг модуль тус бүрээр нь дараах байдлаар тодорхойлов.

\subsection{Хэрэглэгчийн удирдлага}
\begin{enumerate}
\item Хэрэглэгч бүртгүүлэх, нэвтрэх, өөрийн мэдээллээ харах боломжтой байх
\item Профайл зураг, display name зэрэг хувийн мэдээллээ шинэчлэх боломжтой байх
\item Нууц үг солих боломжтой байх
\item JWT authentication ашиглан хамгаалагдсан API-уудыг зөвхөн нэвтэрсэн хэрэглэгч ашиглах
\end{enumerate}

\subsection{Сургалтын агуулгын удирдлага}
\begin{enumerate}
\item Learning track, course, хичээл гэсэн шаталсан бүтэцтэй байх
\item Course нь түвшний ангилалтай байх
\item хичээл нь текст, видео, PDF, хавсралт файл, дадлага ажил мэдээлэл агуулж байх
\item Хэрэглэгч liked хичээлүүд жагсаалтад хичээл хадгалах боломжтой байх
\end{enumerate}

\subsection{Шаталсан сургалтын логик}
\begin{enumerate}
\item Түвшин тогтоох шалгалт-ээр эхний түвшнийг тогтоох
\item Түвшин ахих шалгалт-ээр дараагийн шатанд ахих нөхцөл бүрдүүлэх
\item Course list дээр хэрэглэгчийн түвшинд тохирсон шүүлт хийх
\item Recommendation endpoint нь хэрэглэгчийн түвшин, дууссан хичээл-д тулгуурлан дараагийн content санал болгох
\end{enumerate}

\subsection{Хиймэл оюун дэмжлэгтэй сургалт}
\begin{enumerate}
\item хичээл summary автоматаар үүсгэх
\item хичээл асуулт автоматаар үүсгэх
\item дадлага ажил байхгүй хичээл-д даалгавар үүсгэх
\item Chatbot ашиглан хэрэглэгчид тайлбар өгөх
\item Community пост, comment дээр moderation хийх
\end{enumerate}

\subsection{Ахиц хяналт}
\begin{enumerate}
\item хичээл дуусгалт бүрийг хадгалах
\item Оноог нэгтгэн тооцоолж хэрэглэгчийн нийт оноог шинэчлэх
\item Дуусгасан хичээлийн тоог тооцоолох
\item Долоо хоногийн идэвхийг өдрөөр нь задлан харуулах
\item Курс бүрийн нийт хичээл, дууссан хичээл, ахицын хувь, курсын оноог тооцоолох
\end{enumerate}

\subsection{Харилцаа ба идэвхжүүлэлт}
\begin{enumerate}
\item Community post, comment үүсгэх
\item Зохисгүй агуулгыг moderation-р хянах
\item Games хэсгээр интерактив контент үзүүлэх
\item Landing page дээр системийн танилцуулга, сэтгэгдэл, онцлогуудыг харуулах
\end{enumerate}

\section{Функцийн бус шаардлага}
Системийн чанар, өргөтгөх чадвар, аюулгүй байдал, ашиглахад хялбар байдлыг хангах үүднээс дараах функцийн бус шаардлагыг дэвшүүлэв.

\subsection{Архитектурын шаардлага}
\begin{itemize}
\item Frontend болон backend тусдаа модультай байх
\item REST API дээр суурилсан харилцаатай байх
\item Кодын бүтэц нь app/module тусгаарлалтыг дэмжих
\item Шинэ модуль, endpoint, page нэмэхэд хялбар байх
\end{itemize}

\subsection{Гүйцэтгэлийн шаардлага}
\begin{itemize}
\item Ерөнхий API хариу хурдан байх
\item Course болон хичээл жагсаалтад \texttt{select\_related}, \texttt{prefetch\_related} ашиглаж өгөгдлийн асуулгыг оновчлох
\item Хиймэл оюун endpoint удаашралтай байж болох тул fallback логиктой байх
\item Code runner ашиглах үед хугацаа, санах ой, гаралтын хэмжээнд хязгаар тогтоосон байх
\end{itemize}

\subsection{Аюулгүй байдлын шаардлага}
\begin{itemize}
\item JWT authentication ашиглах
\item CORS болон CSRF trusted origins тохируулгатай байх
\item Production орчинд secure cookie, HTTPS redirect, HSTS тохируулгатай байх
\item Community текст дээр moderation хийх
\item Код ажиллуулах хэсэгт нөөцийн хязгаарлалт тавих
\end{itemize}

\subsection{Найдвартай байдлын шаардлага}
\begin{itemize}
\item Хиймэл оюун үйлчилгээ алдаа өгсөн ч системийн үндсэн урсгал тасалдахгүй байх
\item Database ажиллахгүй бол backend шууд эхлэхээс сэргийлсэн health check ба wait script-тэй байх
\item Debug болон production орчныг тусгаарласан тохиргоотой байх
\end{itemize}

\subsection{Хэрэглэгчийн туршлагын шаардлага}
\begin{itemize}
\item Хичээл, progress, асуулт, discussion-ийг нэг хуудасны орчинд хослуулсан байх
\item Хуудас хоорондын шилжилт ойлгомжтой байх
\item Хэрэглэгчийн одоогийн түвшин, ахиц тодорхой харагддаг байх
\item Хичээл үзэх явцад тусдаа хуудас солихгүйгээр summary, асуулт, дадлага, discussion ашиглах боломжтой байх
\end{itemize}

\section{Use case задлал}
Системийн үндсэн use case-уудыг дараах байдлаар тодорхойлж болно.
\insertimagecustom{zurag-diagram/use-case.png}{Системийн үндсэн хэрэглээний тохиолдлуудын диаграм}{height=0.82\textheight,keepaspectratio}

\subsection{UC1: Бүртгүүлэх ба нэвтрэх}
Оролцогч: Суралцагч.

Урьдчилсан нөхцөл: Хэрэглэгч системд бүртгэлгүй эсвэл бүртгэлтэй байна.

Үндсэн урсгал:
\begin{enumerate}
\item Хэрэглэгч бүртгэлийн форм бөглөж илгээнэ.
\item Систем хэрэглэгчийн мэдээллийг шалгаж, шинэ бүртгэл үүсгэнэ.
\item Хэрэглэгч нэвтрэхэд JWT token авна.
\item Token ашиглан хамгаалагдсан хуудсууд руу орно.
\end{enumerate}

\subsection{UC2: Түвшин тогтоох шалгалт өгөх}
Оролцогч: Суралцагч.

Үндсэн урсгал:
\begin{enumerate}
\item Хэрэглэгч Түвшин тогтоох шалгалт хуудсыг нээнэ.
\item Систем эхний түвшний агуулгад үндэслэн асуултууд үүсгэнэ.
\item Хэрэглэгч хариултаа илгээнэ.
\item Систем оноог тооцож skill level-ийг хадгална.
\item Дараагийн course санал болголт шинэ түвшинд тулгуурлан өөрчлөгдөнө.
\end{enumerate}

\subsection{UC3: хичээл суралцах}
Оролцогч: Суралцагч.

Үндсэн урсгал:
\begin{enumerate}
\item Хэрэглэгч course дотроос хичээл сонгоно.
\item Систем хичээл мэдээлэл болон дадлага ажил-ийг ачаална.
\item Хэрэглэгч үндсэн агуулгыг уншина, видео үзнэ, summary гаргана.
\item асуулт өгч, дадлага ажил ажиллана.
\item хичээл completion хийснээр progress шинэчлэгдэнэ.
\end{enumerate}

\subsection{UC4: Community-д оролцох}
Оролцогч: Суралцагч.

Үндсэн урсгал:
\begin{enumerate}
\item Хэрэглэгч хичээл-тэй холбоотой пост эсвэл comment бичнэ.
\item Систем keyword moderation болон Хиймэл оюун зохицуулалт moderation хийнэ.
\item Зөвшөөрөгдсөн контент нийтлэгдэнэ.
\item Эрсдэлтэй агуулгад warning эсвэл reject шийдвэр гарна.
\end{enumerate}

\section{Мэдээллийн загварын шинжилгээ}
Системийн өгөгдлийн үндсэн объектууд нь харилцан хамааралтай бөгөөд adaptive learning логикийн суурь болдог.
\insertimagecustom{zurag-diagram/er-diagram.png}{Системийн өгөгдлийн бүтцийн entity-relationship диаграм}{height=0.82\textheight,keepaspectratio}

\begin{itemize}
\item \textbf{User}: role, display name, avatar, skill level, Түвшин тогтоох шалгалт төлөв, нийт оноо, completed хичээл, warning count зэрэг талбартай.
\item \textbf{LearningTrack}: ижил чиглэлийн course-уудыг бүлэглэнэ.
\item \textbf{Course}: title, description, level, thumbnail, track холбоос агуулна.
\item \textbf{Lesson}: course-д харьяалагдах ба content, video, file, дадлага ажил, оноо зэрэг мэдээлэл хадгална.
\item \textbf{LessonProgress}: user болон хичээл-ийн хоорондын гүйцэтгэлийн мэдээлэл хадгална.
\item \textbf{LikedLesson}: хэрэглэгчийн сонирхсон хичээл-үүдийг хадгална.
\item \textbf{CommunityPost, CommunityComment}: хэрэглэгчийн харилцааны агуулга хадгална.
\item \textbf{LandingPageContent, LandingPageReview}: системийн танилцуулгын хэсгийн өгөгдөл хадгална.
\end{itemize}

Эдгээр объектуудын холбоо нь системийн бизнес логикийг дэмжинэ. Жишээлбэл хэрэглэгч хичээл дуусгахад \texttt{LessonProgress} шинэчлэгдэж, түүнээс \texttt{User.total\_score}, \texttt{User.completed\_хичээлүүд}, \texttt{User.skill\_level} тооцогдоно. Мөн \texttt{Course.level} болон \texttt{User.skill\_level} талбарууд нь recommendation болон course filtering-д ашиглагдана.

\section{Алгоритмын түвшний шаардлага}
\subsection{Түвшин шинэчлэх логик}
Кодын хэрэгжилтэд нийт оноонд суурилан түвшин шинэчлэх энгийн боловч ойлгомжтой дүрэм ашигласан. Онооны интервалыг ашиглан beginner, elementary, intermediate, advanced түвшинг тодорхойлно. Энэ нь rule-based шийдэл бөгөөд цаашид илүү олон үзүүлэлт оролцуулсан model руу өргөтгөх боломжтой.

\subsection{Recommendation логик}
Recommendation endpoint нь дараах алхмаар ажиллана.
\insertimagecustom{zurag-diagram/flowchart.png}{Түвшин тогтоох шалгалт болон recommendation логикийн ерөнхий урсгал}{height=0.78\textheight,keepaspectratio}
\begin{enumerate}
\item Хэрэглэгчийн дуусгасан хичээлүүдийн ID жагсаалтыг авна.
\item хичээл жагсаалтаас эдгээрийг хасна.
\item Хэрэглэгчийн түвшинд харгалзах курсийн хичээлүүдийг шүүнэ.
\item Курсын гарчиг болон хичээлийн дарааллаар эрэмбэлж эхний 5-г буцаана.
\end{enumerate}

\subsection{Ахицын нэгтгэсэн тооцоолол}
Хичээл дуусгалт бүрийн дараа:
\begin{enumerate}
\item Тухайн хэрэглэгчийн дуусгасан хичээлийн тоо дахин тооцогдоно.
\item Бүх progress мөрийн оноо нийлбэрээр нийт оноо гарна.
\item Нийт оноонд тулгуурлан түвшинг шинэчилнэ.
\item Шаардлага хангасан үед mentor эрх рүү шилжинэ.
\end{enumerate}

\section{Хязгаарлалт ба эрсдэлийн шинжилгээ}
Системийн шинжилгээний явцад дараах хязгаарлалтууд тодорхой болов.
\begin{itemize}
\item Хиймэл оюун -аар үүсгэсэн агуулгын чанар бүрэн тогтвортой биш байж болно.
\item Түвшин тогтоох шалгалт болон Түвшин ахих шалгалт нь content quality-аас шууд хамаарна.
\item Recommendation логик одоогоор rule-based тул хэрэглэгчийн нарийн зан төлөвийг бүрэн тусгахгүй.
\item Community moderation нь keyword болон Хиймэл оюун  шийдлийг хослуулсан ч алдаа гарах магадлалтай.
\item Code runner нь зөвхөн хязгаарлагдмал нөхцөлтэй бөгөөд бүх төрлийн кодыг бүрэн дэмжихгүй.
\end{itemize}

Эдгээр хязгаарлалтыг тайлангийн дараагийн бүлгүүдэд техникийн шийдэл болон цаашдын сайжруулалтын өнцгөөс авч үзнэ.

\section{Бүлгийн дүгнэлт}
Энэ бүлэгт системийн хэрэгцээ, хэрэглэгчийн бүлэг, бизнес процесс, функциональ болон функцийн бус шаардлага, use case, өгөгдлийн загвар, алгоритмын түвшний шаардлагыг тодорхойлов. Шинжилгээний үр дүнд adaptive learning, Хиймэл оюуны тусламж, progress tracking, community interaction, gamification-ийг нэгтгэсэн веб платформ боловсруулах бодит шаардлага батлагдсан. Мөн кодын одоогийн архитектур нь эдгээр шаардлагыг хэрэгжүүлэхүйц модульчлагдсан бүтэцтэй болох нь харагдаж байна.
